\chapter{Transcatheter Aortic Valve Replacement}

\section*{Introduction \& Clinical Indications}
Transcatheter aortic valve replacement (TAVR), also called transcatheter aortic valve implantation, replaces a diseased aortic valve without open-heart valve excision. It is primarily used for severe symptomatic aortic stenosis after evaluation by a multidisciplinary valve team. Age, life expectancy, surgical risk, anatomy, durability considerations, comorbidities, and patient preference influence the choice between TAVR and surgical aortic valve replacement.

TAVR is not appropriate for every valve lesion or every patient. Bicuspid anatomy, active infection, inadequate vascular access, aortic-root anatomy, associated coronary or aortic disease, and the need for other cardiac surgery may favor an alternative strategy. Assessment should be individualized rather than based on age alone.

\section*{Relevant Surgical Anatomy}
The aortic valve lies between the left ventricle and ascending aorta, adjacent to the coronary ostia, mitral-aortic continuity, membranous septum, and conduction system. Transfemoral delivery travels through the iliac arteries and aorta. Valve size and positioning must account for annulus dimensions, leaflet calcification, coronary height, sinus anatomy, and risk of left-ventricular-outflow obstruction.

\section*{Preoperative Workup \& Preparation}
Confirm severity and symptoms with echocardiography and assess ventricular function. ECG, coronary assessment, CT angiography for annular and vascular measurements, renal function, frailty, pulmonary status, dental or infectious issues, and surgical risk are reviewed. Discuss vascular injury, bleeding, stroke, death, paravalvular leak, coronary obstruction, kidney injury, valve failure, pacemaker requirement, and possible conversion to surgery. Plan antithrombotic treatment and access-site management.

\section*{Surgical Technique \& Procedure}
Most procedures use transfemoral arterial access with conscious sedation or general anesthesia depending on patient and center. A guidewire crosses the native valve, balloon preparation may be performed, and the crimped bioprosthetic valve is positioned and expanded during controlled pacing or deployment. Angiography and echocardiography assess valve position, gradients, coronary flow, and regurgitation. Access vessels are closed percutaneously or surgically, and alternative transaxillary, transcarotid, transcaval, or other routes may be selected when femoral access is unsuitable.

\section*{Complications \& Risks}
Important risks include death, stroke, vascular dissection or rupture, bleeding, annular rupture, cardiac tamponade, coronary obstruction, acute kidney injury, conduction block requiring pacemaker, paravalvular leak, valve embolization, infection, and repeat intervention. New neurologic symptoms, chest pain, hypotension, limb ischemia, or loss of pulses requires immediate evaluation.

\section*{Postoperative Care \& Recovery}
Monitor rhythm, conduction, vascular sites, hemoglobin, renal function, neurologic status, and valve performance. Early mobilization is common after uncomplicated transfemoral TAVR. Discharge planning includes medication reconciliation, endocarditis prevention advice, warning signs, echocardiographic follow-up, and coordination with cardiology and primary care.

\section*{Outcomes \& Prognosis}
TAVR can improve survival and symptoms in appropriately selected severe aortic-stenosis patients. Long-term care includes surveillance for structural valve deterioration, thrombosis, endocarditis, conduction problems, and coronary access issues. The heart team should reassess new symptoms rather than attributing them automatically to recovery.

\section*{References}
\begin{enumerate}
\item American College of Cardiology/American Heart Association. Guideline for the Management of Patients With Valvular Heart Disease.
\item European Society of Cardiology/European Association for Cardio-Thoracic Surgery. Guidelines for valvular heart disease.
\item Valve Academic Research Consortium. Standardized definitions for transcatheter valve clinical endpoints.
\end{enumerate}
